%% fig_smith_dynamics.tex
%% Smith ダイナミクスの位相図（2 経路例）
%%
%% 同じ 2 経路モデル（OD 交通量 q=6，c_1=10+f_1，c_2=23-3f_1）。
%% Smith ダイナミクスの ODE：
%%   ḟ_1 = f_2 [c_2 - c_1]_+  -  f_1 [c_1 - c_2]_+
%%        = (6-f_1)(13-4f_1)_+  -  f_1(4f_1-13)_+
%%
%%  f_1 < 3.25: ḟ_1 = (6-f_1)(13-4f_1) > 0  （経路 2 → 経路 1 への乗り換え）
%%  f_1 > 3.25: ḟ_1 = -f_1(4f_1-13)    < 0  （経路 1 → 経路 2 への乗り換え）
%%  f_1 = 3.25: ḟ_1 = 0                       （Nash均衡 = UE）
%%
%%  x 軸：f_1 ∈ [0, 6]，y 軸：ḟ_1 ∈ (-75, 65)
%%
%% Usage: %% fig_smith_dynamics.tex
%% Smith ダイナミクスの位相図（2 経路例）
%%
%% 同じ 2 経路モデル（OD 交通量 q=6，c_1=10+f_1，c_2=23-3f_1）。
%% Smith ダイナミクスの ODE：
%%   ḟ_1 = f_2 [c_2 - c_1]_+  -  f_1 [c_1 - c_2]_+
%%        = (6-f_1)(13-4f_1)_+  -  f_1(4f_1-13)_+
%%
%%  f_1 < 3.25: ḟ_1 = (6-f_1)(13-4f_1) > 0  （経路 2 → 経路 1 への乗り換え）
%%  f_1 > 3.25: ḟ_1 = -f_1(4f_1-13)    < 0  （経路 1 → 経路 2 への乗り換え）
%%  f_1 = 3.25: ḟ_1 = 0                       （Nash均衡 = UE）
%%
%%  x 軸：f_1 ∈ [0, 6]，y 軸：ḟ_1 ∈ (-75, 65)
%%
%% Usage: %% fig_smith_dynamics.tex
%% Smith ダイナミクスの位相図（2 経路例）
%%
%% 同じ 2 経路モデル（OD 交通量 q=6，c_1=10+f_1，c_2=23-3f_1）。
%% Smith ダイナミクスの ODE：
%%   ḟ_1 = f_2 [c_2 - c_1]_+  -  f_1 [c_1 - c_2]_+
%%        = (6-f_1)(13-4f_1)_+  -  f_1(4f_1-13)_+
%%
%%  f_1 < 3.25: ḟ_1 = (6-f_1)(13-4f_1) > 0  （経路 2 → 経路 1 への乗り換え）
%%  f_1 > 3.25: ḟ_1 = -f_1(4f_1-13)    < 0  （経路 1 → 経路 2 への乗り換え）
%%  f_1 = 3.25: ḟ_1 = 0                       （Nash均衡 = UE）
%%
%%  x 軸：f_1 ∈ [0, 6]，y 軸：ḟ_1 ∈ (-75, 65)
%%
%% Usage: %% fig_smith_dynamics.tex
%% Smith ダイナミクスの位相図（2 経路例）
%%
%% 同じ 2 経路モデル（OD 交通量 q=6，c_1=10+f_1，c_2=23-3f_1）。
%% Smith ダイナミクスの ODE：
%%   ḟ_1 = f_2 [c_2 - c_1]_+  -  f_1 [c_1 - c_2]_+
%%        = (6-f_1)(13-4f_1)_+  -  f_1(4f_1-13)_+
%%
%%  f_1 < 3.25: ḟ_1 = (6-f_1)(13-4f_1) > 0  （経路 2 → 経路 1 への乗り換え）
%%  f_1 > 3.25: ḟ_1 = -f_1(4f_1-13)    < 0  （経路 1 → 経路 2 への乗り換え）
%%  f_1 = 3.25: ḟ_1 = 0                       （Nash均衡 = UE）
%%
%%  x 軸：f_1 ∈ [0, 6]，y 軸：ḟ_1 ∈ (-75, 65)
%%
%% Usage: \input{assets/assignment/fig_smith_dynamics.tex}

\begin{tikzpicture}[
  every node/.style={font=\small},
  xscale=1.7,
  yscale=0.030
]

%% ===== 軸 =====
\draw[->] (-0.15, 0) -- (6.5, 0) node[right] {$f_1$};
\draw[->] (0, -70)   -- (0, 82)  node[above] {$\dot{f}_1$};

%% ===== Smith ダイナミクス dF1/dt =====
%% f_1 ∈ [0, 3.25]：ḟ_1 = (6-f_1)(13-4f_1)
\draw[very thick, blue!75!black, samples=34, domain=0:3.25,
      variable=\f]
  plot ({\f}, {(6-\f)*(13-4*\f)});

%% f_1 ∈ [3.25, 6]：ḟ_1 = -f_1*(4*f_1-13)
\draw[very thick, blue!75!black, samples=29, domain=3.25:6,
      variable=\f]
  plot ({\f}, {-\f*(4*\f-13)});

%% ===== UE 平衡点 f_1^* = 3.25 =====
\fill[red!70!black] (3.25, 0) circle [radius=2.2pt];
\node[above right, font=\small, red!70!black] at (3.25, 2)
  {UE\;$f_1^{*}$};

%% ===== 方向矢印（ḟ_1 > 0 → f_1 増加方向，ḟ_1 < 0 → f_1 減少方向）=====
%% 左側（f_1 = 1.5, ḟ_1 = 4.5*7 = 31.5）
\draw[->, thick, green!55!black]
  (1.5, 0) -- (1.9, 0)
  node[above, font=\scriptsize, green!55!black] {$\dot{f}_1>0$};

%% 右側（f_1 = 5.0, ḟ_1 = -5*7 = -35）
\draw[->, thick, green!55!black]
  (5.0, 0) -- (4.6, 0)
  node[above, font=\scriptsize, green!55!black] {$\dot{f}_1<0$};

%% ===== 特徴点のラベル =====
%% 最大値付近：f_1=0 での ḟ_1 = 6*13 = 78
\node[above left, font=\scriptsize] at (0, 78)
  {$f_2 (c_2 - c_1)_+$};

%% f_1=6 での ḟ_1 = -6*(24-13) = -6*11 = -66
\node[below right, font=\scriptsize] at (6, -66)
  {$-f_1(c_1-c_2)_+$};

%% ===== 軸目盛り =====
\node[below, font=\small] at (3.25, 0) {$f_1^{*}$};
\node[below, font=\small] at (0, 0)    {$0$};
\node[below, font=\small] at (6, 0)    {$q$};
\draw[thin, gray] (-0.05, 30) -- (0.05, 30)
  node[left, font=\scriptsize] {$30$};
\draw[thin, gray] (-0.05, 60) -- (0.05, 60)
  node[left, font=\scriptsize] {$60$};
\draw[thin, gray] (-0.05, -30) -- (0.05, -30)
  node[left, font=\scriptsize] {$-30$};
\draw[thin, gray] (-0.05, -60) -- (0.05, -60)
  node[left, font=\scriptsize] {$-60$};

%% ===== ゼロ線（目立つように） =====
\draw[thin, gray!60] (0, 0) -- (6.3, 0);

\end{tikzpicture}


\begin{tikzpicture}[
  every node/.style={font=\small},
  xscale=1.7,
  yscale=0.030
]

%% ===== 軸 =====
\draw[->] (-0.15, 0) -- (6.5, 0) node[right] {$f_1$};
\draw[->] (0, -70)   -- (0, 82)  node[above] {$\dot{f}_1$};

%% ===== Smith ダイナミクス dF1/dt =====
%% f_1 ∈ [0, 3.25]：ḟ_1 = (6-f_1)(13-4f_1)
\draw[very thick, blue!75!black, samples=34, domain=0:3.25,
      variable=\f]
  plot ({\f}, {(6-\f)*(13-4*\f)});

%% f_1 ∈ [3.25, 6]：ḟ_1 = -f_1*(4*f_1-13)
\draw[very thick, blue!75!black, samples=29, domain=3.25:6,
      variable=\f]
  plot ({\f}, {-\f*(4*\f-13)});

%% ===== UE 平衡点 f_1^* = 3.25 =====
\fill[red!70!black] (3.25, 0) circle [radius=2.2pt];
\node[above right, font=\small, red!70!black] at (3.25, 2)
  {UE\;$f_1^{*}$};

%% ===== 方向矢印（ḟ_1 > 0 → f_1 増加方向，ḟ_1 < 0 → f_1 減少方向）=====
%% 左側（f_1 = 1.5, ḟ_1 = 4.5*7 = 31.5）
\draw[->, thick, green!55!black]
  (1.5, 0) -- (1.9, 0)
  node[above, font=\scriptsize, green!55!black] {$\dot{f}_1>0$};

%% 右側（f_1 = 5.0, ḟ_1 = -5*7 = -35）
\draw[->, thick, green!55!black]
  (5.0, 0) -- (4.6, 0)
  node[above, font=\scriptsize, green!55!black] {$\dot{f}_1<0$};

%% ===== 特徴点のラベル =====
%% 最大値付近：f_1=0 での ḟ_1 = 6*13 = 78
\node[above left, font=\scriptsize] at (0, 78)
  {$f_2 (c_2 - c_1)_+$};

%% f_1=6 での ḟ_1 = -6*(24-13) = -6*11 = -66
\node[below right, font=\scriptsize] at (6, -66)
  {$-f_1(c_1-c_2)_+$};

%% ===== 軸目盛り =====
\node[below, font=\small] at (3.25, 0) {$f_1^{*}$};
\node[below, font=\small] at (0, 0)    {$0$};
\node[below, font=\small] at (6, 0)    {$q$};
\draw[thin, gray] (-0.05, 30) -- (0.05, 30)
  node[left, font=\scriptsize] {$30$};
\draw[thin, gray] (-0.05, 60) -- (0.05, 60)
  node[left, font=\scriptsize] {$60$};
\draw[thin, gray] (-0.05, -30) -- (0.05, -30)
  node[left, font=\scriptsize] {$-30$};
\draw[thin, gray] (-0.05, -60) -- (0.05, -60)
  node[left, font=\scriptsize] {$-60$};

%% ===== ゼロ線（目立つように） =====
\draw[thin, gray!60] (0, 0) -- (6.3, 0);

\end{tikzpicture}


\begin{tikzpicture}[
  every node/.style={font=\small},
  xscale=1.7,
  yscale=0.030
]

%% ===== 軸 =====
\draw[->] (-0.15, 0) -- (6.5, 0) node[right] {$f_1$};
\draw[->] (0, -70)   -- (0, 82)  node[above] {$\dot{f}_1$};

%% ===== Smith ダイナミクス dF1/dt =====
%% f_1 ∈ [0, 3.25]：ḟ_1 = (6-f_1)(13-4f_1)
\draw[very thick, blue!75!black, samples=34, domain=0:3.25,
      variable=\f]
  plot ({\f}, {(6-\f)*(13-4*\f)});

%% f_1 ∈ [3.25, 6]：ḟ_1 = -f_1*(4*f_1-13)
\draw[very thick, blue!75!black, samples=29, domain=3.25:6,
      variable=\f]
  plot ({\f}, {-\f*(4*\f-13)});

%% ===== UE 平衡点 f_1^* = 3.25 =====
\fill[red!70!black] (3.25, 0) circle [radius=2.2pt];
\node[above right, font=\small, red!70!black] at (3.25, 2)
  {UE\;$f_1^{*}$};

%% ===== 方向矢印（ḟ_1 > 0 → f_1 増加方向，ḟ_1 < 0 → f_1 減少方向）=====
%% 左側（f_1 = 1.5, ḟ_1 = 4.5*7 = 31.5）
\draw[->, thick, green!55!black]
  (1.5, 0) -- (1.9, 0)
  node[above, font=\scriptsize, green!55!black] {$\dot{f}_1>0$};

%% 右側（f_1 = 5.0, ḟ_1 = -5*7 = -35）
\draw[->, thick, green!55!black]
  (5.0, 0) -- (4.6, 0)
  node[above, font=\scriptsize, green!55!black] {$\dot{f}_1<0$};

%% ===== 特徴点のラベル =====
%% 最大値付近：f_1=0 での ḟ_1 = 6*13 = 78
\node[above left, font=\scriptsize] at (0, 78)
  {$f_2 (c_2 - c_1)_+$};

%% f_1=6 での ḟ_1 = -6*(24-13) = -6*11 = -66
\node[below right, font=\scriptsize] at (6, -66)
  {$-f_1(c_1-c_2)_+$};

%% ===== 軸目盛り =====
\node[below, font=\small] at (3.25, 0) {$f_1^{*}$};
\node[below, font=\small] at (0, 0)    {$0$};
\node[below, font=\small] at (6, 0)    {$q$};
\draw[thin, gray] (-0.05, 30) -- (0.05, 30)
  node[left, font=\scriptsize] {$30$};
\draw[thin, gray] (-0.05, 60) -- (0.05, 60)
  node[left, font=\scriptsize] {$60$};
\draw[thin, gray] (-0.05, -30) -- (0.05, -30)
  node[left, font=\scriptsize] {$-30$};
\draw[thin, gray] (-0.05, -60) -- (0.05, -60)
  node[left, font=\scriptsize] {$-60$};

%% ===== ゼロ線（目立つように） =====
\draw[thin, gray!60] (0, 0) -- (6.3, 0);

\end{tikzpicture}


\begin{tikzpicture}[
  every node/.style={font=\small},
  xscale=1.7,
  yscale=0.030
]

%% ===== 軸 =====
\draw[->] (-0.15, 0) -- (6.5, 0) node[right] {$f_1$};
\draw[->] (0, -70)   -- (0, 82)  node[above] {$\dot{f}_1$};

%% ===== Smith ダイナミクス dF1/dt =====
%% f_1 ∈ [0, 3.25]：ḟ_1 = (6-f_1)(13-4f_1)
\draw[very thick, blue!75!black, samples=34, domain=0:3.25,
      variable=\f]
  plot ({\f}, {(6-\f)*(13-4*\f)});

%% f_1 ∈ [3.25, 6]：ḟ_1 = -f_1*(4*f_1-13)
\draw[very thick, blue!75!black, samples=29, domain=3.25:6,
      variable=\f]
  plot ({\f}, {-\f*(4*\f-13)});

%% ===== UE 平衡点 f_1^* = 3.25 =====
\fill[red!70!black] (3.25, 0) circle [radius=2.2pt];
\node[above right, font=\small, red!70!black] at (3.25, 2)
  {UE\;$f_1^{*}$};

%% ===== 方向矢印（ḟ_1 > 0 → f_1 増加方向，ḟ_1 < 0 → f_1 減少方向）=====
%% 左側（f_1 = 1.5, ḟ_1 = 4.5*7 = 31.5）
\draw[->, thick, green!55!black]
  (1.5, 0) -- (1.9, 0)
  node[above, font=\scriptsize, green!55!black] {$\dot{f}_1>0$};

%% 右側（f_1 = 5.0, ḟ_1 = -5*7 = -35）
\draw[->, thick, green!55!black]
  (5.0, 0) -- (4.6, 0)
  node[above, font=\scriptsize, green!55!black] {$\dot{f}_1<0$};

%% ===== 特徴点のラベル =====
%% 最大値付近：f_1=0 での ḟ_1 = 6*13 = 78
\node[above left, font=\scriptsize] at (0, 78)
  {$f_2 (c_2 - c_1)_+$};

%% f_1=6 での ḟ_1 = -6*(24-13) = -6*11 = -66
\node[below right, font=\scriptsize] at (6, -66)
  {$-f_1(c_1-c_2)_+$};

%% ===== 軸目盛り =====
\node[below, font=\small] at (3.25, 0) {$f_1^{*}$};
\node[below, font=\small] at (0, 0)    {$0$};
\node[below, font=\small] at (6, 0)    {$q$};
\draw[thin, gray] (-0.05, 30) -- (0.05, 30)
  node[left, font=\scriptsize] {$30$};
\draw[thin, gray] (-0.05, 60) -- (0.05, 60)
  node[left, font=\scriptsize] {$60$};
\draw[thin, gray] (-0.05, -30) -- (0.05, -30)
  node[left, font=\scriptsize] {$-30$};
\draw[thin, gray] (-0.05, -60) -- (0.05, -60)
  node[left, font=\scriptsize] {$-60$};

%% ===== ゼロ線（目立つように） =====
\draw[thin, gray!60] (0, 0) -- (6.3, 0);

\end{tikzpicture}
